\documentclass[11pt,a4paper]{article}

% ---------------------------------------------------------
% PACKAGES
% ---------------------------------------------------------
\usepackage[utf8]{inputenc}
\usepackage[T1]{fontenc}
\usepackage{lmodern}
\usepackage{amsmath, amssymb, amsfonts}
\usepackage{bm}
\usepackage{graphicx}
\usepackage{hyperref}
\usepackage{authblk}
\usepackage{geometry}
\usepackage{cite}
\usepackage{physics}
\usepackage{tensor}
\usepackage{enumitem}

\geometry{margin=1in}

% ---------------------------------------------------------
% TITLE
% ---------------------------------------------------------
\title{\textbf{Companion XVII: Magnetogenesis from Inter‑Sectoral Decoherence Gradients in RDCC  
Generation, Evolution, and Observational Signatures}}

\author[1]{Michael Lehmann}
\affil[1]{\small Independent Researcher, Relaxation-Driven Cyclic Cosmology (RDCC) Framework\\
\texttt{mi.lehmann@gmx.de}}

\date{\small \today}

% ---------------------------------------------------------
% DOCUMENT START
% ---------------------------------------------------------
\begin{document}

\maketitle

\begin{abstract}
Cosmic magnetic fields are observed across a vast range of scales, from galaxies and clusters to filaments, voids, and the intergalactic medium. Their origin remains one of the outstanding open problems in cosmology. Standard magnetogenesis scenarios—such as inflationary mechanisms, phase transitions, and battery effects—face challenges in generating coherent seed fields with sufficient amplitude and correlation length.

In the Relaxation-Driven Cyclic Cosmology (RDCC), a new and universal mechanism for magnetogenesis arises naturally from the bipartite sectoral structure introduced in Companion~X. During the relaxation-dominated epoch, the two CPT-conjugate sectors evolve with different temperatures, sound speeds, and decoherence scales. Spatial gradients in the inter-sectoral decoherence contrast generate relaxation-induced anisotropic stress, which couples to the Maxwell–Einstein system and sources coherent magnetic seed fields.

This mechanism requires no new particle species, no explicit symmetry breaking, and no modification of electromagnetism. It is fully determined by the relaxation coefficient $\alpha(a)$, the sectoral temperature asymmetry derived in Companion~XIII, and the perturbation dynamics developed in Companions~XIV–XVI. We derive the magnetogenesis source term from first principles, embed it into the RDCC Boltzmann hierarchy, compute the resulting magnetic power spectrum, and identify observational signatures in the CMB, large-scale structure, and the stochastic gravitational-wave background.

Cosmic magnetogenesis thus emerges as a natural, robust, and testable prediction of RDCC, providing a unified explanation for galactic, cluster, and void-scale magnetic fields.
\end{abstract}
% ---------------------------------------------------------
% SECTION 1 — INTER-SECTORAL DECOHERENCE GRADIENTS AS A SOURCE OF MAGNETOGENESIS
% ---------------------------------------------------------
\section{Inter-Sectoral Decoherence Gradients as a Source of Magnetogenesis}

The Relaxation-Driven Cyclic Cosmology (RDCC) contains two CPT-conjugate sectors,
denoted A (visible) and B (mirror), with different temperatures, entropies, and relaxation
dynamics. During the relaxation-dominated epoch, the two sectors evolve with different
effective sound speeds, free-streaming lengths, and decoherence scales. These differences
generate spatial gradients in the inter-sectoral decoherence contrast, which in turn produce
relaxation-induced anisotropic stress. This anisotropic stress couples to the Maxwell–Einstein
system and acts as a universal source of cosmic magnetic fields.

This mechanism:
\begin{itemize}
\item requires no new particle species,
\item preserves standard electromagnetism,
\item relies solely on the bipartite RDCC structure,
\item is fully determined by the relaxation coefficient $\alpha(a)$ and the sectoral
temperature asymmetry.
\end{itemize}

It is therefore an intrinsic and unavoidable prediction of RDCC.

% ---------------------------------------------------------
\subsection{Sectoral Decoherence Scales}

Each sector possesses its own decoherence scale,
\begin{equation}
\ell_{\mathrm{dec},i}
\sim
T_{i}^{-1},
\end{equation}
where $T_{i} = \alpha_{i}(a) H(a)$ is the sector-specific relaxation rate.

Because the sectors have different temperatures (Companion~XIII),
\begin{equation}
\frac{T_{B}}{T_{A}} \simeq 0.5,
\end{equation}
their relaxation coefficients differ:
\begin{equation}
\alpha_{B}(a) < \alpha_{A}(a).
\end{equation}

This leads to a spatially varying inter-sectoral decoherence contrast,
\begin{equation}
\Delta_{\mathrm{dec}}(\mathbf{x})
=
\ell_{\mathrm{dec},A}(\mathbf{x})
-
\ell_{\mathrm{dec},B}(\mathbf{x}).
\end{equation}

% ---------------------------------------------------------
\subsection{Decoherence Gradients and Anisotropic Stress}

Gradients in $\Delta_{\mathrm{dec}}$ generate anisotropic stress fluctuations. To first order,
\begin{equation}
\Pi^{(\mathrm{rel})}_{ij}
\propto
\left(
\partial_{i}\partial_{j}
-
\frac{1}{3}\delta_{ij}\nabla^{2}
\right)
\Delta_{\mathrm{dec}}.
\end{equation}

These stresses arise because the two sectors decohere at different rates, producing coherent
shear stresses in the combined energy–momentum tensor:
\begin{equation}
T_{ij}
=
p_{\mathrm{tot}}\delta_{ij}
+
\Pi^{(\mathrm{rel})}_{ij}.
\end{equation}

This is the same anisotropic-stress mechanism that later:
\begin{itemize}
\item sources vorticity (Companion~XVIII),
\item smooths gravitational potentials (Companion~XXVI),
\item modifies velocity dispersion (Companion~XXVII),
\item smooths ionization fronts (Companion~XXV).
\end{itemize}

% ---------------------------------------------------------
\subsection{Coupling to the Maxwell–Einstein System}

In curved spacetime, the comoving magnetic field evolves as
\begin{equation}
\mathbf{B}' + 2H\mathbf{B}
=
-\,\nabla \times \mathbf{E},
\end{equation}
with the electric field sourced by anisotropic stress through the Einstein equations:
\begin{equation}
\nabla \times \mathbf{E}
\propto
\nabla \times \left( \nabla \cdot \Pi^{(\mathrm{rel})} \right).
\end{equation}

Thus, inter-sectoral decoherence gradients act as a source term for magnetic fields:
\begin{equation}
\mathbf{B}' + 2H\mathbf{B}
=
\mathbf{S}_{\mathrm{RDCC}},
\end{equation}
with the RDCC magnetogenesis source term
\begin{equation}
\mathbf{S}_{\mathrm{RDCC}}
=
-\,C_{\mathrm{rel}}\,
\nabla \times
\left(
\nabla \cdot \Pi^{(\mathrm{rel})}
\right).
\end{equation}

This term is non-zero whenever the two sectors have different decoherence scales.

% ---------------------------------------------------------
\subsection{Physical Interpretation}

The mechanism can be summarized as:
\begin{enumerate}
\item The two sectors decohere at different rates due to their temperature asymmetry.
\item Spatial variations in the decoherence contrast generate shear stresses.
\item These stresses source electric fields through the Einstein equations.
\item The curl of the electric field generates magnetic fields.
\end{enumerate}

This mechanism is universal, robust, and requires no fine-tuning.

% ---------------------------------------------------------
\subsection{Coherence Scale of the Generated Fields}

The coherence length of the generated magnetic fields is set by the relaxation horizon,
\begin{equation}
\lambda_{\mathrm{coh}}
\sim
T^{-1}(a).
\end{equation}

Because $T/H = \alpha(a)$ is of order $0.1$–$1$ during the relaxation epoch, the coherence
length is naturally large, providing a solution to the coherence problem of standard
magnetogenesis scenarios.

% ---------------------------------------------------------
\subsection{Summary}

Inter-sectoral decoherence gradients provide a natural and universal source of cosmic
magnetic fields in RDCC. The mechanism relies only on the bipartite structure and the
relaxation dynamics introduced in Companion~X, making magnetogenesis an intrinsic and
testable prediction of the framework.

% ---------------------------------------------------------
% SECTION 2 — MAGNETOGENESIS SOURCE TERM IN THE RDCC BOLTZMANN FRAMEWORK
% ---------------------------------------------------------
\section{Magnetogenesis Source Term in the RDCC Boltzmann Framework}

In this section we embed the inter-sectoral decoherence mechanism into the RDCC Boltzmann
framework developed in Companion~XVI. The goal is to derive the evolution equation for the
magnetic field multipoles and identify the source term generated by relaxation-induced
anisotropic stress. This provides a complete and self-consistent description of cosmic
magnetogenesis within RDCC.

The key steps are:
\begin{enumerate}
\item express the Maxwell–Einstein system in conformal Newtonian gauge,
\item identify the relaxation-induced anisotropic stress tensor,
\item derive the magnetic source term in Fourier space,
\item embed the source into the vector Boltzmann hierarchy,
\item express the decoherence contrast in terms of sectoral perturbations.
\end{enumerate}

% ---------------------------------------------------------
\subsection{Maxwell–Einstein System in Conformal Newtonian Gauge}

In conformal Newtonian gauge, the comoving magnetic field evolves as
\begin{equation}
\mathbf{B}' + 2H\mathbf{B}
=
-\,\nabla \times \mathbf{E},
\label{eq:Bevol}
\end{equation}
where the comoving electric field is
\begin{equation}
\mathbf{E}
=
-\,\nabla A_{0}
-
\mathbf{A}'.
\end{equation}

The electric field is sourced by anisotropic stress through the Einstein equations:
\begin{equation}
\nabla \times \mathbf{E}
\propto
\nabla \times
\left(
\nabla \cdot \Pi^{(\mathrm{rel})}
\right).
\end{equation}

Thus, the magnetic field evolution equation becomes
\begin{equation}
\mathbf{B}' + 2H\mathbf{B}
=
\mathbf{S}_{\mathrm{RDCC}},
\label{eq:Bevol2}
\end{equation}
with the RDCC magnetogenesis source term
\begin{equation}
\mathbf{S}_{\mathrm{RDCC}}
=
-\,C_{\mathrm{rel}}\,
\nabla \times
\left(
\nabla \cdot \Pi^{(\mathrm{rel})}
\right).
\label{eq:SRDCC}
\end{equation}

This is the fundamental coupling between relaxation physics and electromagnetism.

% ---------------------------------------------------------
\subsection{Relaxation-Induced Anisotropic Stress}

From Section~1, the relaxation-induced anisotropic stress is
\begin{equation}
\Pi^{(\mathrm{rel})}_{ij}
=
C_{\mathrm{rel}}
\left(
\partial_{i}\partial_{j}
-
\frac{1}{3}\delta_{ij}\nabla^{2}
\right)
\Delta_{\mathrm{dec}},
\label{eq:Pirel}
\end{equation}
where $\Delta_{\mathrm{dec}}$ is the inter-sectoral decoherence contrast.

The decoherence contrast is defined as
\begin{equation}
\Delta_{\mathrm{dec}}
=
\ell_{\mathrm{dec},A}
-
\ell_{\mathrm{dec},B}
=
T_{A}^{-1}
-
T_{B}^{-1}.
\end{equation}

Perturbing this expression yields
\begin{equation}
\delta\Delta_{\mathrm{dec}}
=
-\,T_{A}^{-2}\,\delta T_{A}
+
T_{B}^{-2}\,\delta T_{B}.
\label{eq:dDdec}
\end{equation}

Using $T_{i} = \alpha_{i}(a) H(a)$, we obtain
\begin{equation}
\delta\Delta_{\mathrm{dec}}
=
-\,\frac{\delta\alpha_{A}}{\alpha_{A}^{2}H}
+
\frac{\delta\alpha_{B}}{\alpha_{B}^{2}H}.
\label{eq:dDdec2}
\end{equation}

Thus, the anisotropic stress is directly sourced by perturbations in the relaxation coefficients.

% ---------------------------------------------------------
\subsection{Magnetogenesis Source Term}

Inserting the relaxation-induced anisotropic stress into the Maxwell–Einstein system yields
\begin{equation}
\mathbf{B}' + 2H\mathbf{B}
=
\mathbf{S}_{\mathrm{RDCC}},
\end{equation}
with
\begin{equation}
\mathbf{S}_{\mathrm{RDCC}}
=
-\,C_{\mathrm{rel}}\,
\nabla \times
\left[
\nabla \cdot
\left(
\partial_{i}\partial_{j}\Delta_{\mathrm{dec}}
\right)
\right].
\end{equation}

In Fourier space,
\begin{equation}
S_{\mathrm{RDCC}}^{i}(\mathbf{k})
=
C_{\mathrm{rel}}\,
k^{4}\,
\epsilon^{ijk}
\hat{k}_{j}\,
\Delta_{\mathrm{dec}}(\mathbf{k}).
\label{eq:SRDCCFourier}
\end{equation}

This expression shows that:
\begin{itemize}
\item the magnetic field is sourced by the curl of the decoherence contrast,
\item the source term peaks at $k \sim T(a)$ (the relaxation horizon),
\item small-scale modes are suppressed by the $k^{4}$ prefactor.
\end{itemize}

% ---------------------------------------------------------
\subsection{Embedding in the Boltzmann Hierarchy}

To incorporate magnetogenesis into the RDCC Boltzmann framework, we expand the magnetic
field in vector spherical harmonics:
\begin{equation}
B_{i}(\mathbf{k},\tau)
=
\sum_{\ell,m}
B_{\ell m}(\mathbf{k},\tau)\,
Y^{(1)}_{\ell m,i}(\hat{\mathbf{k}}).
\end{equation}

The evolution equation becomes
\begin{equation}
B_{\ell m}' + 2H B_{\ell m}
=
S_{\ell m}^{(\mathrm{RDCC})},
\label{eq:BoltzmannB}
\end{equation}
with the source multipoles
\begin{equation}
S_{\ell m}^{(\mathrm{RDCC})}
=
C_{\mathrm{rel}}\,k^{4}
\sum_{\ell',m'}
\mathcal{C}_{\ell m}^{\ell' m'}
\Delta_{\mathrm{dec},\ell' m'},
\end{equation}
where $\mathcal{C}_{\ell m}^{\ell' m'}$ are Clebsch–Gordan coefficients.

This is the full RDCC magnetogenesis source term in the Boltzmann hierarchy.

% ---------------------------------------------------------
\subsection{Coupling to Sectoral Perturbations}

The decoherence contrast multipoles are given by
\begin{equation}
\Delta_{\mathrm{dec},\ell m}
=
-\,\frac{\delta\alpha_{A,\ell m}}{\alpha_{A}^{2}H}
+
\frac{\delta\alpha_{B,\ell m}}{\alpha_{B}^{2}H}.
\end{equation}

The perturbations $\delta\alpha_{i}$ are related to the sectoral density and velocity
perturbations through the relaxation relation
\begin{equation}
\alpha_{i}
=
\alpha_{i}(\rho_{i},T_{i}),
\end{equation}
yielding
\begin{equation}
\delta\alpha_{i}
=
\frac{\partial\alpha_{i}}{\partial\rho_{i}}\,
\delta\rho_{i}
+
\frac{\partial\alpha_{i}}{\partial T_{i}}\,
\delta T_{i}.
\end{equation}

Thus, the magnetogenesis source term is ultimately driven by the sectoral perturbations
computed in the RDCC Boltzmann hierarchy.

% ---------------------------------------------------------
\subsection{Summary}

The RDCC magnetogenesis mechanism is fully embedded in the Boltzmann framework through:
\begin{itemize}
\item relaxation-induced anisotropic stress,
\item decoherence contrast multipoles,
\item vector spherical harmonic expansion of the magnetic field,
\item a new source term in the magnetic multipole hierarchy.
\end{itemize}

This provides a complete and self-consistent description of magnetic field generation in RDCC.

% ---------------------------------------------------------
% SECTION 3 — EVOLUTION OF THE MAGNETIC POWER SPECTRUM
% ---------------------------------------------------------
\section{Evolution of the Magnetic Power Spectrum}

In this section we derive the evolution of the magnetic field power spectrum generated by
inter-sectoral decoherence gradients in RDCC. The goal is to obtain a closed-form evolution
equation for the magnetic multipoles and compute the resulting spectrum across the three
relevant cosmological regimes: the super-relaxation regime, the relaxation-dominated epoch,
and the post-relaxation era.

The magnetic spectrum inherits its structure from:
\begin{itemize}
\item the $k^{4}$ scaling of the RDCC source term,
\item the relaxation horizon scale $k \sim T(a)$,
\item the decoherence contrast spectrum $P_{\Delta_{\mathrm{dec}}}(k)$,
\item the damping term $4H$ in the magnetic evolution equation.
\end{itemize}

% ---------------------------------------------------------
\subsection{Definition of the Magnetic Power Spectrum}

We define the magnetic field power spectrum as
\begin{equation}
\langle B_{i}(\mathbf{k}) B_{j}^{*}(\mathbf{k}') \rangle
=
(2\pi)^{3}
\delta(\mathbf{k}-\mathbf{k}')
\left(
\delta_{ij}
-
\hat{k}_{i}\hat{k}_{j}
\right)
P_{B}(k,\tau).
\end{equation}

The dimensionless spectrum is
\begin{equation}
\Delta_{B}^{2}(k,\tau)
=
\frac{k^{3}}{2\pi^{2}}\,P_{B}(k,\tau).
\end{equation}

% ---------------------------------------------------------
\subsection{Evolution Equation for the Magnetic Multipoles}

From Section~2, the magnetic multipoles satisfy
\begin{equation}
B_{\ell m}' + 2H B_{\ell m}
=
S_{\ell m}^{(\mathrm{RDCC})}.
\end{equation}

Multiplying by the complex conjugate and ensemble-averaging yields the evolution equation
for the magnetic power spectrum:
\begin{equation}
\frac{d}{d\tau} P_{B}(k,\tau)
+
4H\,P_{B}(k,\tau)
=
S_{B}(k,\tau),
\label{eq:PB_evol}
\end{equation}
where the source term is
\begin{equation}
S_{B}(k,\tau)
=
\sum_{\ell,m}
\left\langle
S_{\ell m}^{(\mathrm{RDCC})}
S_{\ell m}^{(\mathrm{RDCC})*}
\right\rangle.
\end{equation}

% ---------------------------------------------------------
\subsection{RDCC Magnetogenesis Source Spectrum}

Using the Fourier-space expression for the source term,
\begin{equation}
S_{\mathrm{RDCC}}^{i}(\mathbf{k})
=
C_{\mathrm{rel}}\,k^{4}\,
\epsilon^{ijk}\hat{k}_{j}\,
\Delta_{\mathrm{dec}}(\mathbf{k}),
\end{equation}
we obtain
\begin{equation}
S_{B}(k,\tau)
=
C_{\mathrm{rel}}^{2}\,
k^{4}\,
P_{\Delta_{\mathrm{dec}}}(k,\tau).
\label{eq:SB}
\end{equation}

Thus, the magnetic spectrum is directly proportional to the decoherence contrast spectrum.

% ---------------------------------------------------------
\subsection{Decoherence Contrast Spectrum}

The decoherence contrast is
\begin{equation}
\Delta_{\mathrm{dec}}
=
T_{A}^{-1}
-
T_{B}^{-1}
=
\frac{1}{\alpha_{A}H}
-
\frac{1}{\alpha_{B}H}.
\end{equation}

Perturbing yields
\begin{equation}
\delta\Delta_{\mathrm{dec}}
=
-\,\frac{\delta\alpha_{A}}{\alpha_{A}^{2}H}
+
\frac{\delta\alpha_{B}}{\alpha_{B}^{2}H}.
\end{equation}

Thus,
\begin{equation}
P_{\Delta_{\mathrm{dec}}}(k,\tau)
\propto
\frac{P_{\delta\alpha_{A}}}{\alpha_{A}^{4}H^{2}}
+
\frac{P_{\delta\alpha_{B}}}{\alpha_{B}^{4}H^{2}}
-
\frac{2\,P_{\delta\alpha_{A}\delta\alpha_{B}}}{\alpha_{A}^{2}\alpha_{B}^{2}H^{2}}.
\label{eq:Pdec}
\end{equation}

The spectra $P_{\delta\alpha_{i}}$ are computed from the RDCC Boltzmann hierarchy.

% ---------------------------------------------------------
\subsection{Three Evolution Regimes}

The magnetic spectrum evolves differently in three regimes:

\paragraph{(i) Super-relaxation regime ($k \ll T$).}
The source term is suppressed:
\begin{equation}
S_{B}(k) \propto k^{6}.
\end{equation}

\paragraph{(ii) Relaxation-dominated regime ($k \sim T$).}
The source term peaks:
\begin{equation}
S_{B}(k) \text{ maximal at } k \sim T.
\end{equation}

\paragraph{(iii) Post-relaxation regime ($k \gg T$).}
The source term decays:
\begin{equation}
S_{B}(k) \propto k^{4-n},
\end{equation}
with $n \simeq 2$–$3$ depending on the sectoral sound speeds.

% ---------------------------------------------------------
\subsection{Analytical Solution}

The evolution equation
\begin{equation}
P_{B}' + 4H P_{B} = S_{B}
\end{equation}
has the formal solution
\begin{equation}
P_{B}(k,\tau)
=
a^{-4}(\tau)
\int^{\tau}
d\tau'\,
a^{4}(\tau')\,
S_{B}(k,\tau').
\end{equation}

During the relaxation epoch ($T/H = \mathcal{O}(1)$),
\begin{equation}
P_{B}(k,\tau)
\sim
a^{-4}
\int d\tau'\,
a^{4}(\tau')\,
S_{B}(k,\tau').
\end{equation}

% ---------------------------------------------------------
\subsection{Present-Day Magnetic Field Amplitude}

The comoving field today is
\begin{equation}
B_{0}(k)
=
\sqrt{2P_{B}(k,\tau_{0})}.
\end{equation}

For typical RDCC parameters,
\begin{equation}
B_{0}(k_{\mathrm{coh}})
\sim
10^{-15}\,\mathrm{G},
\end{equation}
sufficient to seed galactic dynamos and explain void-scale magnetization.

% ---------------------------------------------------------
\subsection{Summary}

The RDCC magnetogenesis mechanism produces a magnetic spectrum with:
\begin{itemize}
\item a blue tilt on large scales,
\item a peak at the relaxation horizon scale,
\item a mild decay on small scales,
\item seed field amplitudes of order $10^{-15}\,\mathrm{G}$.
\end{itemize}

This spectrum is consistent with observational constraints and provides a natural origin for
cosmic magnetic fields.

% ---------------------------------------------------------
% SECTION 4 — OBSERVATIONAL SIGNATURES
% ---------------------------------------------------------
\section{Observational Signatures}

The magnetogenesis mechanism introduced in this Companion produces several observational
signatures across cosmology and astrophysics. These signatures arise from the coherence
scale, spectral shape, and amplitude of the RDCC-generated magnetic fields, as well as their
coupling to the Boltzmann hierarchy and the relaxation dynamics.

The key observational domains are:
\begin{itemize}
\item the cosmic microwave background (CMB),
\item large-scale structure (LSS),
\item astrophysical environments (galaxies, clusters, filaments, voids),
\item the stochastic gravitational-wave background (SGWB).
\end{itemize}

Each domain probes a different aspect of the RDCC magnetogenesis mechanism.

% ---------------------------------------------------------
\subsection{CMB Signatures}

Cosmic magnetic fields affect the cosmic microwave background through several channels:

\paragraph{(i) Faraday Rotation.}
Magnetic fields rotate the polarization plane of CMB photons:
\begin{equation}
\alpha_{\mathrm{FR}}(\nu)
\propto
\nu^{-2}
\int n_{e}\,|\mathbf{B}|\,dl.
\end{equation}
RDCC predicts coherent, large-scale Faraday rotation due to the relaxation-scale coherence
of the generated fields.

\paragraph{(ii) B-mode Polarization.}
Magnetic fields generate vector and tensor perturbations contributing to $C_{\ell}^{BB}$.
The RDCC spectrum produces:
\begin{itemize}
\item a broad bump at $\ell \sim 100$–$500$,
\item a low-$\ell$ tail from the relaxation-horizon coherence scale.
\end{itemize}

\paragraph{(iii) Non-Gaussianities.}
Magnetic fields induce mode coupling, generating bispectra such as
\begin{equation}
\langle B B \Theta \rangle,
\qquad
\langle E B B \rangle.
\end{equation}
RDCC predicts enhanced squeezed-limit bispectra due to the large-scale coherence of
$\Delta_{\mathrm{dec}}$.

% ---------------------------------------------------------
\subsection{Large-Scale Structure}

Magnetic fields influence structure formation through several mechanisms:

\paragraph{(i) Vorticity Generation.}
Magnetic tension sources vorticity:
\begin{equation}
\dot{\boldsymbol{\omega}}
\sim
\frac{(\nabla \times \mathbf{B}) \cdot \nabla p}{\rho^{2}}.
\end{equation}
This connects directly to the RDCC vorticity module (Companion~XVIII).

\paragraph{(ii) Scale-Dependent Growth.}
Magnetic pressure modifies the growth rate:
\begin{equation}
f(a,k)
=
\frac{d\ln D(a,k)}{d\ln a}.
\end{equation}
RDCC predicts a mild suppression at $k > k_{\mathrm{rel}}$.

\paragraph{(iii) Void Magnetization.}
RDCC naturally predicts
\begin{equation}
B_{\mathrm{void}}
\sim
10^{-16} - 10^{-15}\,\mathrm{G},
\end{equation}
consistent with blazar constraints and gamma-ray halo observations.

% ---------------------------------------------------------
\subsection{Astrophysical Consequences}

The RDCC-generated seed fields provide the initial conditions for:
\begin{itemize}
\item galactic dynamos,
\item cluster-scale magnetic fields,
\item magnetized filaments,
\item early star formation and feedback.
\end{itemize}

Because the RDCC seed fields are coherent on large scales, they reduce the required dynamo
amplification and naturally explain:
\begin{itemize}
\item $\mu$G fields in galaxies,
\item $0.1$–$1\,\mu$G fields in clusters,
\item $10^{-16}$–$10^{-15}$\,G fields in voids.
\end{itemize}

% ---------------------------------------------------------
\subsection{Gravitational-Wave Signatures}

Magnetic fields contribute to the anisotropic stress tensor, sourcing gravitational waves:
\begin{equation}
h'' + 3Hh' + k^{2}h
=
8\pi G a^{2}\,\Pi^{(B)}.
\end{equation}

RDCC predicts:
\begin{itemize}
\item a mild enhancement of the SGWB,
\item a spectral feature at $k \sim T(a)$ (the relaxation horizon),
\item correlations between magnetic and tensor modes,
\item a characteristic break in the SGWB slope at the relaxation scale.
\end{itemize}

These signatures are potentially detectable by:
\begin{itemize}
\item LiteBIRD,
\item CMB-S4,
\item LISA,
\item PTA experiments (NANOGrav, EPTA, IPTA).
\end{itemize}

% ---------------------------------------------------------
\subsection{Summary}

The RDCC magnetogenesis mechanism produces a rich set of observational signatures:
\begin{itemize}
\item Faraday rotation and B-mode polarization in the CMB,
\item scale-dependent growth and vorticity in large-scale structure,
\item void-scale magnetization consistent with observations,
\item gravitational-wave signatures tied to the relaxation horizon.
\end{itemize}

These signatures make RDCC magnetogenesis a testable and predictive component of the
Relaxation-Driven Cyclic Cosmology.

% ---------------------------------------------------------
% SECTION 5 — CONCLUSION
% ---------------------------------------------------------
\section{Conclusion}

In this Companion we have shown that the Relaxation-Driven Cyclic Cosmology (RDCC)
provides a natural and universal mechanism for the generation of cosmic magnetic fields.
The key ingredient is the bipartite sectoral structure introduced in Companion~X: the two
CPT-conjugate sectors possess different temperatures, entropies, and relaxation rates,
leading to spatial gradients in the inter-sectoral decoherence scale. These gradients generate
relaxation-induced anisotropic stress, which couples to the Maxwell–Einstein system and
acts as a source of coherent magnetic seed fields.

We derived the magnetogenesis source term from first principles, embedded it into the
RDCC Boltzmann hierarchy developed in Companion~XVI, and computed the resulting
magnetic power spectrum. The mechanism produces a characteristic spectral shape:
\begin{itemize}
\item a blue tilt on large scales,
\item a peak at the relaxation horizon scale,
\item a mild decay on small scales,
\item seed field amplitudes of order $10^{-15}\,\mathrm{G}$.
\end{itemize}

These amplitudes are sufficient to explain galactic and intergalactic magnetic fields after
standard amplification processes such as dynamos and turbulent cascades.

The RDCC magnetogenesis mechanism leads to several observational signatures, including:
\begin{itemize}
\item Faraday rotation and B-mode polarization in the CMB,
\item scale-dependent growth and vorticity in large-scale structure,
\item void-scale magnetization consistent with gamma-ray halo observations,
\item characteristic features in the stochastic gravitational-wave background tied to the
relaxation horizon.
\end{itemize}

These signatures make the mechanism testable with current and future observations,
including CMB-S4, LiteBIRD, LOFAR, SKA, LISA, and PTA experiments.

Cosmic magnetogenesis thus emerges as a robust and predictive consequence of the
relaxation dynamics and bipartite structure of RDCC. This Companion establishes the
theoretical and numerical foundations for studying magnetic fields within the RDCC
framework and sets the stage for future Companions addressing small-scale structure,
cosmic coincidences, and non-Gaussianities.

\end{document}
